\reviewercomment{6. Mathematical model

The gLV model successfully reproduces the qualitative transition in coalescence outcomes and serves as a useful minimal framework. However, it is limited to purely competitive interactions and does not explicitly capture environmentally mediated effects such as pH modification.

We therefore suggest framing the model more explicitly as a phenomenological rather than mechanistic description of the system.

In addition, while the robustness analysis across interaction coefficient distributions is appreciated, a brief discussion of the biological interpretation of these distributions would be helpful. In particular, clarifying how the mean interaction strength parameter ($\mu$) relates to underlying ecological processes would improve interpretability.}

\responselabel We agree that the gLV model should be presented as a phenomenological competitive baseline, not as a complete mechanistic model of the experimental system. We have therefore revised the model description to state explicitly that the gLV framework explores how the statistical properties of effective interaction coefficients shape coalescence outcomes, while omitting explicit environmental state variables such as pH, metabolite concentrations, and carrying-capacity differences. In the model, $\alpha_{ij}$ is now described as an effective per-capita inhibitory term for the effect of species $j$ on species $i$; it can absorb the net steady-state consequences of direct competition and indirect inhibition, including pH-mediated inhibition, but it does not identify the biochemical mechanism and does not dynamically model pH itself. Because we restrict $\alpha_{ij} \geq 0$, we also now state that the simulations represent competitive or growth-inhibiting interactions only, and that facilitative interactions such as cross-feeding are not represented.

We also clarified the biological interpretation of the interaction coefficient distributions and of $\mu$. The uniform, Gaussian, and Gamma coefficient distributions are not intended as fitted empirical distributions of molecular mechanisms. They are phenomenological ensembles of heterogeneous effective competitive effects, parameterized so that their mean and variance are matched; the robustness analysis asks whether the Mixture-to-Dominance/Restructuring transition depends on the particular distributional shape. In that context, $\mu$ is an effective interaction-intensity parameter: within the simulated gLV system, it is the mean of the off-diagonal competition coefficient distribution and therefore controls average competitive suppression. In the experiments, we do not treat nutrient concentration as a direct measurement of $\mu$. Instead, nutrient enrichment is interpreted as increasing net invasion resistance, measured by failed pairwise invasions among the 12 most abundant isolates (Nutr$-$: $2 \pm 1\%$; Base: $33 \pm 4\%$; Nutr$+$: $48 \pm 4\%$), with several possible underlying ecological processes, including denser resource competition, increased metabolic activity, environmental modification such as pH shifts, and changes in the competition-facilitation balance.

This reframing preserves the main use of the model while making its limits explicit. The competition-only gLV model shows that competitive assembly structure is sufficient to generate frequent Dominance and positive within-parental-community selection correlation, but it does not claim that competition is the only mechanism operating in the experiments. Indeed, the revised Results and Discussion now emphasize that pH modification by dominant taxa provides a mechanistic explanation for winner identity in the Nutr$+$ top-down regime: pH mismatch did not significantly increase Dominance frequency within Base or Nutr$+$, but among acid-alkaline coalescence pairs the acidic parental community won 38/44 Nutr$+$ events (86.4\%, binomial $p = 9.4 \times 10^{-7}$). Thus, the model is used to identify a sufficient coarse-grained mechanism for community-level selection, whereas the experimental interpretation allows environmental mediation to explain additional features of the data.

\mschange{We revised the Results to call the gLV model a phenomenological framework rather than a mechanistic model of pH modification, metabolic cross-feeding, or carrying-capacity variation. We revised the Supplementary Methods to define $\alpha_{ij}$ as an effective inhibitory coefficient, to state that $\alpha_{ij} \geq 0$ makes the model competitive-only, and to define $\mu$ as the mean intensity of effective competitive suppression. We also revised the Discussion to state that the mapping between nutrient concentration and $\mu$ is approximate and reflects net interaction intensity measured through invasion resistance, not a unique biochemical mechanism. No new response figure was added because the concern is addressed by textual clarification and by existing main-text and supplementary analyses.}
